\section{A distinguished large ball in real dimension six}
\label{six:section}

In complex dimension three, the hyperplane construction of
Proposition~\ref{app:large} would use a plane conic.  Its genus is zero,
so it cannot supply the handle required by Lemma~\ref{surf:packing}.
We instead use a smooth plane cubic to construct a model for the shell
outside a reserved ball.  Throughout this section we use rational parameters
\begin{equation}\label{six:parameters}
 \frac12<a<1,\qquad s=1-a,\qquad
 b=\frac{2-a}{3}=\frac{1+s}{3}.
\end{equation}
In particular \(0<s<b\).  All projective spaces and varieties in this
section are complex.  The notation \(\cO_C(j)\), for a plane curve \(C\),
means the restriction of \(\cO_{\PP^2}(j)\).

The target is the shell
\[
 \{z\in\C^3:a<\pi\textstyle\sum_j|z_j|^2<1\}
\]
of Lemma~\ref{geo:shell}.  We first identify a quadric normal model that
transfers into this shell.  A lower circle cut of an auxiliary threefold
realizes that model.  We then construct a model over a cubic curve and
transfer its compact regions into the auxiliary threefold, preserving
the cutting moment.

\subsection{The quadric normal model}

Let \(T=\Bl_o\PP^3\), where \(o\notin H_\infty\), and put
\(L=H-aE\).  Here \(H\) is the hyperplane pullback and \(E\) is the
exceptional divisor.  Choose a smooth quadric through \(o\), general
with respect to \(H_\infty\), and let \(D\) be its strict transform.

\begin{lemma}\label{six:quadric}
There is an identification \(D=\Bl_{q_1,q_2}\PP^2\), for two distinct
points \(q_1,q_2\). Write \(M\) for the pullback of the plane hyperplane
class and \(e_i\) for the exceptional divisor class over \(q_i\).
Under this identification,
\begin{equation}\label{six:quadric-classes}
 H|_D=2M-e_1-e_2,\qquad E|_D=M-e_1-e_2.
\end{equation}
Consequently
\begin{equation}\label{six:normal-class}
 N_{D/T}=\cO_D(3M-e_1-e_2),\qquad
 L|_D=3bM-s(e_1+e_2).
\end{equation}
For every rational \(0<c<s\), deformation to the normal cone of \(D\)
has ample component classes
\[
 (1-2c)H-(a-c)E
 \quad\hbox{and}\quad
 3bM-s(e_1+e_2)+cU
\]
on its strict component and on
\begin{equation}\label{six:Y}
 Y=\PP_D\bigl(\mathbf1\oplus\cO_D(3M-e_1-e_2)\bigr),
\end{equation}
respectively.  The intersections of \(D\) with \(E\) and \(H_\infty\)
are disjoint smooth curves; their images in \(\PP^2\) are a line
\(\ell\) and a smooth conic \(Q\), both through \(q_1,q_2\).
\end{lemma}

\begin{proof}
Projection from \(o\) on the quadric becomes a morphism after blowing
up \(o\).  It contracts the strict transforms of the two ruling lines
through \(o\), which are disjoint curves of self-intersection \(-1\).
One can see both the morphism and its inverse in the equation
\[
 x_0x_1=x_2x_3,\qquad o=[1:0:0:0].
\]
Projection is to \([x_1:x_2:x_3]\); the inverse is the quadratic map
\[
 [u:v:w]\longmapsto[vw:u^2:uv:uw],
\]
whose two simple base points are \([0:1:0]\) and \([0:0:1]\).
Its resolution identifies \(D\) with the stated two-point blowup.
The inverse quadratic system gives the first equality in
\eqref{six:quadric-classes}; the exceptional curve over \(o\) is the
strict transform of \(u=0\), giving the second equality.
Since \([D]=2H-E\), restriction gives \eqref{six:normal-class}.

The strict component class is ample because
\[
 a-c>0,\qquad (1-2c)-(a-c)=s-c>0.
\]
The restrictions of the class on \(Y\) to its zero and infinity
sections are \(L|_D\) and \((L-cD)|_D\), and its fiber degree is \(c>0\).
These conditions imply ampleness here.  Indeed, subtract a sufficiently
small ample class on \(Y\); the two section restrictions remain ample
and the fiber degree remains positive. By Hirschowitz's invariant-cycle
argument, recalled in \cite[Section~2]{FultonMacPhersonSottileSturmfels1995},
every curve is rationally, hence numerically, equivalent to an effective
cycle invariant under scalar multiplication in the normal line.
Its irreducible components lie in the fixed sections or in fibers.
The perturbed class is therefore nef. Adding back the small ample
class proves ampleness by \cite[Corollary~1.4.10]{Lazarsfeld2004}.

Finally, \(E\) and \(H_\infty\) are disjoint, and their restrictions
to a general \(D\) are smooth.  Their classes in
\eqref{six:quadric-classes} identify their plane images as the line
and conic asserted.  The conic is smooth for a general choice of the
quadric.  Their strict transforms are disjoint; equivalently, their
only plane intersections are the two marked points, with distinct
tangent directions there.
\end{proof}

The normal-cone transfer of Proposition~\ref{geo:transfer} will be
applied to \(Y\) away from infinity and from the two vertical
submanifolds over \(D\cap E\) and \(D\cap H_\infty\).  We next realize
the required part of \(Y\) by a cut of another threefold.

\subsection{The lower cut and its complex structure}

Put
\[
 V_0=\PP_{\PP^2}\bigl(\mathbf1\oplus\cO_{\PP^2}(3)\bigr),
 \qquad
 \widehat V=\Bl_{(q_1,0),(q_2,0)} V_0 .
 \]
The two points lie in the zero section.  Write \(U_1\) for the
pullback of the dual tautological class on \(V_0\), and
\(\widehat E_i\) for the point exceptional divisors.  The circle
rotating the \(\cO(3)\) summand lifts to \(\widehat V\).

The next lemma assumes an invariant K\"ahler representative.
The cubic family below supplies it: Lemma~\ref{six:cubic-ampleness}
establishes its polarization, and the subsequent full-section
construction and circle average give the required form.

\begin{lemma}\label{six:cut}
Suppose \(s<d_1<b\) and \(\widehat V\) has an invariant Kähler form
in the class
\begin{equation}\label{six:Vhat-class}
 A=3bM+d_1U_1-s(\widehat E_1+\widehat E_2).
\end{equation}
Normalize its first moment \(p_1\) to be zero on the strict bottom
section \(D=\Bl_{q_1,q_2}\PP^2\).  For \(0<c<s\), the cut retaining
\(p_1\le c\) is the complex projective bundle \(Y\) of
\eqref{six:Y}, with an invariant Kähler form in class
\[
 3bM-s(e_1+e_2)+cU.
\]
Its open retained region is symplectically identified with
\(\{p_1<c\}\subset\widehat V\), and this identification preserves
the projection to \(D\) on the attracting basin of the bottom
section.  In particular it preserves avoidance of inverse images
of \(\ell\cup Q\) in the plane.
\end{lemma}

\begin{proof}
The circle is Hamiltonian: \(H^1(\widehat V;\R)=0\), since the
projective bundle over \(\PP^2\) has vanishing first cohomology
and point blowups preserve it.  Thus contraction
of the invariant form with its generating vector field is exact.
The bottom and upper sections are fixed.  At either point before
blowing up, the tangent weights are \(0,0,1\), the first two being
the plane directions.  Thus on the exceptional \(\PP^2\) the fixed
loci are the projectivized base directions, a line belonging to
the bottom section, and the pure vertical point.  A line connecting
these loci has area \(s\) and weight one.  Its moment difference
is therefore \(s\).  An ordinary projective fiber has area \(d_1\)
and weight one at the bottom.  It follows that the fixed levels
are exactly
\begin{equation}\label{six:fixed-levels}
 0,\quad s,\quad d_1.
\end{equation}

The attracting basin \(\mathcal B\) of the bottom section for
complex scalar contraction is the total space of
\[
 N'=\cO_D(3M-e_1-e_2).
\]
To verify this including the exceptional directions, use base
coordinates \(y_1,y_2\) at a marked point and an original line
coordinate \(z\).  In the blowup chart with nonzero first base
direction,
\[
 y_2=y_1v,\qquad z=y_1w.
\]
The contraction acts by \(w\mapsto\lambda w\).  On the other base
chart \(w'=w/v\), so this coordinate is a fiber coordinate in
\(\pi^*\cO(3)\otimes\cO_D(-e_1-e_2)\).
The charts with nonzero base direction, together with the original
charts away from the marked points, give its entire total space.
The omitted pure vertical direction does not contract to the
bottom.  This proves the description of \(\mathcal B\).

We emphasize that \(\mathcal B\) is an attracting basin, not a
moment sublevel.  It contains \(\{p_1<s\}\).  Indeed, contraction
decreases the moment strictly along every nonconstant complex
orbit.  Its limit exists on the projective variety and is fixed.
For a point of moment less than \(s\), \eqref{six:fixed-levels}
forces this limit to lie in the bottom section.  Along a nonzero
fiber of \(N'\), the opposite limit is at level \(s\) or \(d_1\).

Form the Kähler reduction of \(\widehat V\times\C\) by the diagonal
circle at
\begin{equation}\label{six:cut-equation}
 p_1+\pi|v|^2=c.
\end{equation}
Every point of this level lies over \(\mathcal B\), since its
first moment is at most \(c<s\).  In
\(N'\oplus\mathbf1\), each nonzero pair \((w,v)\) has a unique
circle in its complex scalar orbit satisfying
\eqref{six:cut-equation}.  At contraction the total moment tends
to zero.  At expansion it tends to infinity if \(v\ne0\), and
to \(s\) or \(d_1\) if \(v=0\); all these limits exceed \(c\).
Strict moment increase proves the uniqueness.  The circle acts
freely there because the fiber weights are one.

The holomorphic quotient of the nonzero pairs is
\(\PP_D(N'\oplus\mathbf1)\).  The preceding orbit description
identifies it with the reduction, with its Kähler complex
structure.  This is also the usual holomorphic-quotient
description of Kähler cutting; see
\cite[Section~2]{BurnsGuilleminLerman2002}.
For completeness, the quotient map is holomorphic and constant
on complex scalar orbits.  The complex orthogonal complement
of an orbit at the level identifies its complex tangent quotient
with the tangent space of this projectivization.  Positivity of
the ambient Kähler form therefore gives precisely the asserted
Kähler structure.

The map
\[
 x\longmapsto
 \bigl[x,\sqrt{(c-p_1(x))/\pi}\,\bigr],\qquad p_1(x)<c,
\]
uses a positive real additional coordinate and is symplectic:
the pullback of its standard area form vanishes.  It preserves
the base point in \(D\), although it is not in general holomorphic.
The quotient class restricts on the bottom to
\(3bM-s(e_1+e_2)\).  Its projective fiber area is \(c\), from the
residual weight-one moment interval \([0,c]\).
These two data determine the class on this projective bundle.
The base-preserving description also proves the last assertion.
\end{proof}

\subsection{A cubic degeneration and the ample class on its component}

Choose a smooth plane cubic \(C\) through \(q_1,q_2\), general
enough to meet \(\ell\cup Q\) in finitely many points and to be
smooth at the marked points.  Such a choice exists in the linear
system of cubics through the two points. Away from the marks, products
of a line vanishing at each mark but not at the test point, times an
arbitrary linear form, separate values and tangent first jets.
Singular vanishing at a fixed unmarked point therefore imposes three
independent linear conditions; the incidence over the two-dimensional
plane has smaller dimension than the section space. At each mark,
vanishing of the differential is a proper linear condition, also
avoided by a general member. One can also avoid the proper subspaces
of cubics containing \(\ell\) or \(Q\).
The curve has genus one
and \(\deg_C M=3\).

Let
\[
 \mathcal Z=\Bl_{C\times\{0\}}(\PP^2\times\mathbb A^1).
 \]
Its central components are the strict plane \(P\) and
\(Z=\PP_C(\mathbf1\oplus\cO_C(3))\).  Denote the dual tautological
class on \(Z\) by \(U_2\), and put
\[
 \mathcal N=\cO_{\mathcal Z}(3M-[Z]).
 \]
Since \([Z]|_Z=-U_2\) and \([Z]|_P=3M\), its restrictions are
\begin{equation}\label{six:N-restrictions}
 \mathcal N|_Z=\cO_Z(3M+U_2),\qquad
 \mathcal N|_P=\mathbf1.
\end{equation}
Take \(\PP_{\mathcal Z}(\mathbf1\oplus\mathcal N)\), and blow it
up along the two parameter sections obtained from the constant
points \(q_i\) in the zero section.  More precisely, use the
strict transforms in \(\mathcal Z\) of
\(\{q_i\}\times\mathbb A^1\), followed by the zero section of
this projective bundle.  Let \(f:\mathcal X\to\mathbb A^1\)
be the resulting family.

These two sections are disjoint and pass through the smooth
submersion locus of the central fiber.  Indeed, in the local
normal-cone blowup of \((x,t)\), with \(x=0\) defining \(C\),
the strict transform of the constant point is \(x/t=0\) in
the \(t\)-chart.  It meets \(Z\) at its zero normal coordinate,
away from \(P\cap Z\); its first projective coordinate is also
zero.  Thus the total space \(\mathcal X\) is smooth and its
central fiber has two smooth components:
\[
 W=\Bl_{(q_1,0,0),(q_2,0,0)}
       \PP_Z(\mathbf1\oplus\cO_Z(3M+U_2)),
 \qquad
 W'=P\times\PP^1.
\]
The double locus is the first projective bundle over the infinity
section of \(Z\); the two blowup centers miss it.  Every nonzero
fiber of \(f\) is \(\widehat V\).  The family is flat: the smooth
integral total space has local rings torsion-free over the local
discrete valuation ring of the parameter curve, and such modules
are flat.

For rational parameters
\begin{equation}\label{six:d-parameters}
 s<d_1<d_2<b
\end{equation}
consider the relative rational class
\begin{equation}\label{six:relative-class}
 \mathcal A=3bM-d_2[Z]+d_1U_1
              -s(\widehat E_1+\widehat E_2).
\end{equation}
Here \(U_1\) is the dual tautological class for the first bundle
in the family, and the exceptional classes now denote the
family blowups.  Its restrictions are
\begin{equation}\label{six:component-classes}
 \mathcal A|_W=3bM+d_2U_2+d_1U_1
                      -s(\widehat E_1+\widehat E_2),
 \qquad
 \mathcal A|_{W'}=3(b-d_2)M+d_1U_1.
\end{equation}

\begin{lemma}\label{six:cubic-ampleness}
Both classes in \eqref{six:component-classes} are ample.
Consequently \(\mathcal A\) is relatively ample near zero.
\end{lemma}

\begin{proof}
The class on \(W'\) is a positive product class.  On \(W\) use
the two algebraic scalar factors, with the second factor lifted
to the first bundle using its linearization on \(U_2\).
An ordinary surface fiber over \(C\) is
\(\PP_{\PP^1}(\mathbf1\oplus\cO_{\PP^1}(1))\).
Its toric diagram is
\begin{equation}\label{six:diagram}
 \mathcal Q=\{p_1,p_2\ge0:\ p_1\le d_1,\ p_1+p_2\le d_2\}.
\end{equation}
The first projective fibers have size \(d_1\); on their two
endpoint sections the second sizes are \(d_2\) and \(d_2-d_1\).
This gives the four vertices
\[
 (0,0),\ (d_1,0),\ (d_1,d_2-d_1),\ (0,d_2).
\]
The invariant boundary curves have positive areas
\(d_1,d_2-d_1,d_1,d_2\).

Over either marked point the fiber consists of the strict
surface fiber and the exceptional plane.  The former has the
corner \((0,0)\) cut by \(p_1+p_2\ge s\); its five edge lengths
are
\[
 s,\quad d_1-s,\quad d_2-d_1,\quad d_1,\quad d_2-s.
\]
The exceptional plane has size \(s\).  All these numbers are
positive by \eqref{six:d-parameters}.  The torus acts torically
on both components: its tangent characters at the blown-up
point are \(0,(1,0),(0,1)\), which induce the standard toric
action on the exceptional plane.

A horizontal invariant irreducible curve must be a strict
vertex section.  To see the exhaustiveness of this assertion,
intersect it with the generic fiber over \(C\).  This is a finite
set preserved by the connected torus, hence fixed pointwise.
The curve is therefore generically in the fixed locus, whose
four horizontal components are precisely the vertex sections.
Their degrees, in the order of the displayed vertices, are
\begin{equation}\label{six:vertex-degrees}
 9b-2s,\qquad 9(b-d_1),\qquad
 9(b-d_2),\qquad 9(b-d_2).
\end{equation}
Indeed \(U_2\) restricts to \(0\) and \(-3M\) on its two
sections.  On the first of them the upper first section has
\(U_1=-3M\), whereas on the second \(\mathcal N\) is trivial
and both first sections have \(U_1=0\).
Only the bottom vertex section passes through the two blowup
centers, losing \(2s\).  All numbers in
\eqref{six:vertex-degrees} are positive; the first is \(3+s\).

We explain why these curve tests prove ampleness.  The invariant
irreducible curves just enumerated have only finitely many
numerical classes.  The boundary curves in the ordinary fibers
vary in algebraic families with constant numerical class, and
there are only two special fibers.  Fix an ample class \(B\)
on \(W\).  For sufficiently small \(\epsilon>0\), every listed
curve has nonnegative degree against
\(\mathcal A|_W-\epsilon B\).
Every curve is rationally, hence numerically, equivalent to a
torus-invariant effective cycle
\cite[Section~2]{FultonMacPhersonSottileSturmfels1995}.
This result applies to a connected solvable group on a projective
variety and does not require finitely many orbits. Its irreducible
components are invariant, since a connected group cannot permute
finitely many components nontrivially. Thus
\(\mathcal A|_W-\epsilon B\) is nef.
The sum of a nef class and a positive ample class is ample
\cite[Corollary~1.4.10]{Lazarsfeld2004},
proving the assertion on \(W\).
Finally, ampleness on the reduced central fiber is equivalent
to ampleness on its components, and relative ampleness is open
near that fiber.  These standard projective ampleness facts
apply to a sufficiently divisible integral multiple of
\(\mathcal A\); see \cite{Hartshorne1977,Lazarsfeld2004,Stacks}.
\end{proof}

Choose a sufficiently divisible \(N\) so that \(N\mathcal A\)
is integral, relatively very ample near zero, and satisfies
relative Serre vanishing for the ideal of \(W\).
Every section on \(W\) then lifts locally in the parameter:
the exact sequence with \(I_W\otimes\cO(N\mathcal A)\) gives
a surjection onto \(H^0(W,\cO(N\mathcal A)|_W)\).
Choose a full basis on \(W\) diagonal by the two torus weights
and lift it.  Add any further sections needed for a relative
embedding, subtracting their restrictions using these lifts.
The added sections vanish on \(W\).  The induced projective
form, divided by \(N\), consequently restricts to the form from
the chosen full diagonal basis on \(W\).
This restricted form is invariant under the compact two-torus,
since its characters preserve the diagonal Hermitian norm.

The first circle acts on the entire family and preserves the
blowup centers.  Average the total form over this circle.
The average is closed and Kähler on the fibers, and it leaves
the prescribed restriction on \(W\) unchanged.  If needed, a
positive form from the parameter disk makes the total form
Kähler without changing any fiber restriction.  In particular,
the nonzero fibers acquire invariant Kähler forms in the class
\eqref{six:Vhat-class}, as required by Lemma~\ref{six:cut}.
No extension of the second torus action to the family is used.

\subsection{Toric coordinates and the integrated base area}

To apply Lemma~\ref{surf:packing}, we need the horizontal area on the
punctured cubic as a function of the two fiber moments.  The full
projective system records the two point blowups as forced vanishing
of its coefficient sections.  We will use these vanishing orders to
compute the area lost when the marked points are removed.

On the affine chart of \(Z\), \(U_2\) is trivialized by the nonvanishing
first homogeneous coordinate.  Thus the two affine line coordinates
\(w_1,w_2\) both have transition bundle \(\cO_C(3)\).

\begin{lemma}\label{six:weights}
For the above full projective system, the weight \((k_1,k_2)\)
runs over
\[
 0\le k_1\le Nd_1,\qquad 0\le k_2\le Nd_2-k_1.
\]
Its coefficient space on \(C\) is
\begin{equation}\label{six:weight-space}
 H^0\!\left(C,\,
 \cO_C((3bN-3k_1-3k_2)M)
       \bigl(-m_kq_1-m_kq_2\bigr)\right),
 \qquad m_k=(sN-k_1-k_2)_+.
\end{equation}
For sufficiently large divisible \(N\), these systems after
removing the indicated zeros are all basepoint-free.
\end{lemma}

\begin{proof}
Expand first in the projective coordinate associated to \(U_1\).
A term of degree \(k_1\) leaves coefficient class
\[
 (3bN-3k_1)M+(Nd_2-k_1)U_2.
\]
Expanding this in the second projective coordinate gives precisely
the range and the coefficient line in \eqref{six:weight-space}.
Sections on the point blowup are exactly sections before the
blowup vanishing to total order at least \(sN\) at each center.
In local coordinates \((y,w_1,w_2)\), the coefficient of
\(w_1^{k_1}w_2^{k_2}\) must therefore vanish in \(y\) to order
at least \(m_k\), with no other condition.

Writing \(t=(k_1+k_2)/N\), the degree of the line in
\eqref{six:weight-space}, divided by \(N\), is
\[
 9(b-t)-2(s-t)_+.
\]
This decreases on \(0\le t\le d_2\), and is at least
\(9(b-d_2)>0\).  Its integral degrees therefore tend to infinity
uniformly over the weights as \(N\) increases through divisible
values.  On a genus-one curve every line of degree at least two
is basepoint-free: for each point, the evaluation map is
surjective by the vanishing of \(H^1\) after subtracting that
point; see \cite[Tags~0E3B and~0E3C]{Stacks}.
This proves the assertion.
\end{proof}

\begin{lemma}\label{six:toric-model}
Let \(\Delta\) be a compact downward polytope strictly below the
two upper boundaries of \(\mathcal Q\), and let
\[
 C^\circ=C\setminus\{q_1,q_2\}.
\]
Over every compact bordered surface
\(\Sigma\subset C^\circ\), the corresponding affine moment
region has smooth product disk coordinates, including their
lower axes, in which its form is
\begin{equation}\label{six:surface-form}
 \Omega=\omega_V+\dd\Gamma,\qquad
 \dd_C\Gamma=\kappa(p)>0.
\end{equation}
Here \(\Gamma\) is horizontal and toric, with smooth extensions
to neighborhoods of \(\Sigma\) and \(\Delta\).  For an exhaustion
of \(C^\circ\) by such surfaces, the integrated areas converge
uniformly on \(\Delta\) to
\begin{equation}\label{six:area}
 H_0(p)=9(b-p_1-p_2)-2(s-p_1-p_2)_+.
\end{equation}
The same limit and uniform convergence hold after deleting any
additional finite set of base points.
\end{lemma}

\begin{proof}
\emph{Moment coordinates and the horizontal form.}
In a local holomorphic frame of \(M|_C\), the projective potential
on the unmarked affine chart is
\begin{equation}\label{six:potential}
 B(y,\rho)=\frac1N\log\left(\sum_k g_k(y)e^{k\cdot\rho}\right),
 \qquad \rho_j=\log|w_j|^2.
\end{equation}
Here \(g_k\) is the sum of squared absolute values of a basis of
the coefficient space in \eqref{six:weight-space}, in its local
frame.  It is positive away from the marked points.  Near a mark,
with local coordinate \(y=0\), it has exactly the form
\begin{equation}\label{six:coefficient-zero}
 g_k(y)=|y|^{2m_k}a_k(y),\qquad a_k>0
\end{equation}
with \(a_k\) smooth up to zero, by basepoint-freeness after
factoring off the required zeros.

The moments are \(p=\partial_\rho B\).
For fixed \(y\), the Hessian \(B_{\rho\rho}\) is a positive
multiple of the covariance matrix of the weight vectors with
positive coefficients, and is positive definite since the
weights span the plane.  Its gradient is a diffeomorphism onto
the interior of \(\mathcal Q\).  Indeed, for \(p\) in that
interior the function \(B-p\cdot\rho\) is proper and strictly
convex: the maximum of the linear forms
\((k/N-p)\cdot\rho\) grows at least linearly in \(|\rho|\).
It has a unique minimum, smoothly depending on \(y,p\).

We check the axes explicitly.  Put \(r_j=|w_j|^2\) and write
\[
 F(y,r)=N^{-1}\log\sum_k g_k(y)r_1^{k_1}r_2^{k_2}.
\]
Then \(p_j=r_jF_{r_j}\).  At \(r_j=0\), \(F_{r_j}>0\),
because the complete system contains weights with \(k_j=1\).
The derivative of \(p_j\) in another radial variable vanishes
on this face.  On its active variables, the logarithmic Hessian
of the face system is positive definite.  Thus the Jacobian
of \(r\mapsto p\) is block triangular with invertible diagonal
blocks at every lower face, including the origin.
The inverse radii are smooth through those faces, and
\[
 z_j=w_j\sqrt{F_{r_j}/\pi},\qquad \pi|z_j|^2=p_j,
\]
are smooth disk coordinates there.
Compactness gives these coordinates on neighborhoods of all
the prescribed truncated regions, with positive margins from
the upper faces and the omitted base points.

For fixed \(p\) set
\begin{equation}\label{six:legendre}
 G(y,p)=\inf_\rho\bigl(B(y,\rho)-p\cdot\rho\bigr).
\end{equation}
This partial Legendre transform is the usual reduced K\"ahler potential;
see \cite[Sections~4--5]{BurnsGuillemin2004}.
On lower faces this is interpreted using the corresponding
face system.  The envelope identity gives
\(\dd_yG=\dd_yB\) at the selected radii.  Therefore the form
in moment and phase coordinates is
\begin{equation}\label{six:local-form}
 \sum_{j=1}^2\dd p_j\wedge\dd\theta_j
       +\dd\bigl(\dd_y^cG\bigr).
\end{equation}
The one-form \(\dd_y^cG\) extends smoothly through the axes:
it is \(\dd_y^cF\) evaluated at the smooth inverse radii.
Possible terms \(p_j\log p_j\) in \(G\) are independent of the
base and disappear upon this differentiation.

The base form \(\kappa(p)=\dd_y\dd_y^cG\) is positive.
For instance, on the open orbit put \(\zeta_j=\log w_j\),
so \(\rho_j=\zeta_j+\overline{\zeta_j}\).
The positive Hermitian Hessian of \(B\) has blocks
\[
 \begin{pmatrix}
 B_{y\bar y}&B_{y\rho}\\
 B_{\rho\bar y}&B_{\rho\rho}
 \end{pmatrix}.
\]
Differentiating \(B_\rho=p\) at fixed \(p\) yields
\[
 G_{y\bar y}=
 B_{y\bar y}-B_{y\rho}B_{\rho\rho}^{-1}B_{\rho\bar y}>0.
\]
At an axis, restrict the original Kähler form to the coordinate
complex submanifold and use the same calculation on the active
variables.  At the origin it is the positive restriction to
the zero section.

These base forms are global.  If another frame of \(M\) is
\(e'=he\), put \(\lambda=\log|h|^2\) and
\(\phi=\arg(h)/(2\pi)\).  The changes are
\[
 \rho'_j=\rho_j-3\lambda,\quad
 \theta'_j=\theta_j-3\phi,\quad
 B'=B-3b\lambda,\quad
 G'=G-3(b-p_1-p_2)\lambda.
\]
Since \(\dd^c\lambda=\dd\phi\), the primitive in
\eqref{six:local-form} changes by \(-3b\,\dd\phi\), whose
exterior derivative is zero.
Thus the local potentials have the transition law of the real class
\begin{equation}\label{six:curvature-class}
 3(b-p_1-p_2)c_1(M|_C).
\end{equation}
We will use this class after extending the curvature across the marked
points.

Over a bordered \(\Sigma\), choose a smooth frame of \(M\),
which is possible since \(H^2(\Sigma;\Z)=0\).
In the resulting product disk coordinates, subtracting the
standard fiber form leaves
\[
 \eta=\kappa(y,p)+
       \sum_j\dd p_j\wedge\alpha_j(y,p),
\]
where the \(\alpha_j\) are smooth base one-forms.
There are no fiber--fiber terms.  Closure implies
\(\partial_{p_j}\kappa=\dd_C\alpha_j\) and
\(\partial_{p_i}\alpha_j=\partial_{p_j}\alpha_i\).
Choose \(\dd_C\gamma_0=\kappa(y,0)\) and define
\[
 \Gamma(y,p)=\gamma_0+
       \int_0^1\sum_j p_j\alpha_j(y,tp)\,\dd t.
\]
Radial homotopy in the downward domain gives
\(\dd\Gamma=\eta\).  All choices can be made on a slightly
larger bordered surface and a slightly larger truncated
polytope, proving \eqref{six:surface-form} with the asserted
neighborhood extensions.

\medskip\noindent
\emph{The two marked-point contributions to the base area.}
The form \eqref{six:surface-form} now gives the required local model.
To determine its limiting integrated area, we compute the curvature
concentrated at the omitted marks.  Write \(\xi=k/N\) and
\(v(\xi)=(s-\xi_1-\xi_2)_+\).  By
\eqref{six:coefficient-zero}, for
\(\lambda=\log|y|^2<0\) the potential \(B\) differs by a
uniformly bounded amount from
\[
 A_\lambda(\rho)=
       \max_{\xi\in\mathcal Q\cap N^{-1}\Z^2}
            \bigl(\xi\cdot\rho+\lambda v(\xi)\bigr).
\]
The bound is uniform in \(\rho\) and \(y\) near the mark:
there are finitely many \(a_k\), bounded above and below by
positive constants.
Take \(N\) divisible enough to include the vertices of the
subdivision of \(\mathcal Q\) by \(\xi_1+\xi_2=s\).
Barycentric interpolation in either cell, on which \(v\) is
affine, shows that for every \(p\in\mathcal Q\) and every
\(\rho\),
\[
 A_\lambda(\rho)-p\cdot\rho\ge\lambda v(p).
\]
This bound is sharp.  If \(p_1+p_2\le s\), take
\(\rho=(\lambda,\lambda)\), for which
\(A_\lambda=\lambda s\).
If \(p_1+p_2\ge s\), take \(\rho=0\), for which
\(A_\lambda=0\).
Taking infima and using the bounded difference proves
\begin{equation}\label{six:log-mass}
 G(y,p)=(s-p_1-p_2)_+\log|y|^2+O(1).
\end{equation}
This argument also applies on axes and at \(p_1+p_2=s\);
indeed its bounded error is uniform in \(p\).

For each fixed \(p\), subtracting the logarithmic term in
\eqref{six:log-mass} leaves a bounded subharmonic function on
the punctured coordinate disk.  It extends subharmonically
over the mark and has no atomic curvature there: an atom
would give an unbounded logarithmic singularity.
With our normalization
\(\dd\dd^c\log|y|^2=\delta_0\), the extended current of \(G\)
has an atom of mass \((s-p_1-p_2)_+\) at each mark.
Its total mass, by \eqref{six:curvature-class}, is
\(9(b-p_1-p_2)\).  Equivalently, compare its local potentials
with smooth potentials in that real class; the difference
is a global locally integrable function and its exact
curvature has total integral zero.  Removing the two atoms
gives precisely \eqref{six:area}.

The reduced form is smooth at every other base point, so
deleting finitely many further points changes none of these
integrals.  Along an increasing bordered exhaustion the
integrated positive forms give continuous functions on
the compact set \(\Delta\), increasing pointwise to the
continuous function \(H_0\).
Dini's theorem makes their convergence uniform.
\end{proof}

\subsection{Moment-preserving transfer and the packing}

We record the normalization issue for the first transfer out
of the cubic component.

\begin{lemma}\label{six:moment-transfer}
Let \(\Sigma\) be a connected compact bordered surface in the unmarked
base, and let \(K\) be the entire compact moment region over \(\Sigma\)
corresponding to a full-dimensional downward polytope \(\Delta\) as
in Lemma~\ref{six:toric-model}, strictly below the two upper faces
of \(\mathcal Q\).
For sufficiently small nonzero parameter, symplectic
parallel transport from \(W\) to \(\widehat V\) is defined
on a neighborhood of \(K\) and preserves \(p_1\) exactly,
when the moments on both components are zero on their
first bottom sections.
If \(K\) avoids the inverse image of a closed subset of
the plane under the family map, its image has the same
avoidance for sufficiently small parameter.
\end{lemma}

\begin{proof}
The region avoids the central double locus, which is the
second upper face \(p_1+p_2=d_2\), and it avoids the two
point blowups.  It is therefore in the submersion locus.
Parallel transport for the averaged total form exists for
short time on a neighborhood of this compact set by
Lemma~\ref{geo:projective-transfer}.  The connection is invariant
under the first circle, so transport \(\Phi\) is both
symplectic and equivariant.  Hence
\[
 \dd(p_1^{\widehat V}\circ\Phi)=\dd p_1^W
\]
there.  The difference is constant on the connected region.
The region includes points with first affine coordinate
zero.  They are fixed points on the first bottom section,
away from the double locus and the blowups.
The corresponding fixed component in the total space is
the strict transform of the global zero section.
Short equivariant transport of these points stays in that
component, so both normalized moments are zero.  The
constant is consequently zero.
Finally, the relevant inverse image is closed in the total
space.  A compact set disjoint from it stays disjoint under
sufficiently short transport.
\end{proof}

\begin{proposition}\label{six:large}
Let \(R_0,R_1,\ldots,R_k>0\), with \(R_0\ge1/2\), satisfy
\[
 \sum_{i=0}^k R_i^3<1,\qquad
 R_i+R_j<1\quad(i\ne j).
\]
Then the disjoint union of their closed six-dimensional
balls embeds symplectically into the open ball of capacity
one, with each embedding defined on a neighborhood of its
closed source.
\end{proposition}

\begin{proof}
If \(k=0\) the assertion is immediate.  Otherwise the
strict inequalities allow a rational \(a>R_0\), with
\(1/2<a<1\), and auxiliary \(r_i>R_i\), \(1\le i\le k\),
such that
\begin{equation}\label{six:auxiliary}
 r_i<s=1-a,\qquad
 \sum_{i=1}^k r_i^3<1-a^3.
\end{equation}
Use the notation \eqref{six:parameters}.
The limiting moment domain is
\[
 \Delta_\infty=
 \{p_1,p_2\ge0:\ p_1\le s,\ p_1+p_2\le b\}.
\]
Since \(b>s\), direct integration of \eqref{six:area} gives
\begin{align}
 \int_{\Delta_\infty}H_0(p)\,\dd p_1\dd p_2
 &=\int_0^s
       \left(\frac92(b-u)^2-(s-u)^2\right)\,\dd u
       \label{six:volume}\\
 &=\frac32\bigl(b^3-(b-s)^3\bigr)-\frac{s^3}{3}
   =\frac{s}{2}-\frac{s^2}{2}+\frac{s^3}{6}
   =\frac{1-a^3}{6}.\notag
\end{align}
The full cap of size \(r_i<s\) is supported inside this
domain and has integral
\[
 \int_{p_1,p_2\ge0}(r_i-p_1-p_2)_+\,
                 \dd p_1\dd p_2=\frac{r_i^3}{6}.
\]

Fix \(d_1\) with \(s<d_1<b\).  Choose rational \(h_1,c,h_2,d_2\)
with
\begin{equation}\label{six:truncation}
 \max_i r_i<h_1<c<s<d_1<d_2<b,\qquad
 s<h_2<d_2,
\end{equation}
so close to \(h_1=s\) and \(h_2=b\) that on
\[
 \Delta=\{p_1,p_2\ge0:\ p_1\le h_1,\ p_1+p_2\le h_2\}
\]
the function
\[
 P(p)=H_0(p)-\sum_{i=1}^k(r_i-p_1-p_2)_+
\]
has positive integral, and hence positive mean.
Figure~\ref{fig:sixdim-moment} shows the truncated and limiting
domains, the cut level, and the change of slope of the limiting density.
\begin{figure}[tbp]
\centering
\begin{tikzpicture}[font=\small,>=Stealth]
 \begin{scope}[x=13cm,y=13cm]
  \fill[black!7] (0,0)--(.21,0)--(.21,.17)--(0,.38)--cycle;
  \draw[dashed,thick] (0,0)--(.3,0)--(.3,.133333)--(0,.433333)--cycle;
  \draw[thick] (0,0)--(.21,0)--(.21,.17)--(0,.38)--cycle;
  \draw[densely dotted,thick] (0,.3)--(.3,0);
  \draw[dash dot] (.26,0)--(.26,.25)
    node[above,align=center] {cut\\\(p_1=c\)};
  \draw[->] (0,0)--(.48,0) node[right] {\(p_1\)};
  \draw[->] (0,0)--(0,.48) node[above] {\(p_2\)};
  \node[below left] at (0,0) {\(0\)};
  \node[left] at (0,.433333) {\(b\)};
  \node[left] at (0,.38) {\(h_2\)};
  \node[left] at (0,.3) {\(s\)};
  \node[below] at (.21,0) {\(h_1\)};
  \node[below] at (.26,0) {\(c\)};
  \node[below] at (.3,0) {\(s\)};
  \node at (.08,.12) {\(\Delta\)};
  \node[above right] at (.075,.358333) {\(\Delta_\infty\)};
  \node[rotate=-45,fill=white,inner sep=1pt] at (.065,.235)
    {\(p_1+p_2=s\)};
 \end{scope}
 \begin{scope}[shift={(8.1cm,.5cm)},x=10cm,y=1.05cm]
  \draw[->] (0,0)--(.48,0) node[right] {\(t\)};
  \draw[->] (0,0)--(0,3.8) node[above] {\(H_0(t)\)};
  \draw[thick] (0,3.3)--(.3,1.2)--(.433333,0);
  \draw[dotted] (.3,0)--(.3,1.2);
  \node[left] at (0,3.3) {\(9b-2s\)};
  \node[below] at (.3,0) {\(s\)};
  \node[below] at (.433333,0) {\(b\)};
  \node[above right] at (.12,2.46) {slope \(-7\)};
  \node[right] at (.35,.8) {slope \(-9\)};
  \node[align=center] at (.23,-1.05)
    {\(t=p_1+p_2\)\\
     \(H_0(t)=9(b-t)-2(s-t)_+\)};
 \end{scope}
\end{tikzpicture}
\caption{The six-dimensional packing domain and its limiting base area.
The solid domain \(\Delta\) lies entirely below the cut \(p_1=c\).
The dashed domain \(\Delta_\infty\) is used to compute the limiting
volume \((1-a^3)/6\). The dotted line marks the change of slope of
\(H_0\); the portion of \(\Delta\) below it remains in the packing domain.
Each of the two marked points of the elliptic base contributes the
subtracted mass \((s-t)_+\). The drawing illustrates
\(h_1<c<s<h_2<b\); the full parameter order is
\eqref{six:truncation}.}
\label{fig:sixdim-moment}
\end{figure}

Such choices are possible by \eqref{six:auxiliary},
\eqref{six:volume}, and convergence of these domains to
\(\Delta_\infty\).  The cap supports are already wholly
inside every sufficiently close truncation.

The hypotheses of Lemma~\ref{surf:packing} hold.
First, \(H_0\) is positive on \(\Delta\).
As a function of \(t=p_1+p_2\), its slopes are \(-7\)
for \(t<s\) and \(-9\) for \(t>s\), so it is concave.
Each negative cap is concave as well, proving concavity
of \(P\).
The auxiliary simplices lie strictly away from the two
outer faces.  Choose \(\tau<1\) with
\(\tau h_1>\max_i r_i\) and \(\tau h_2>\max_i r_i\).
If \(p\notin\tau\Delta\), either
\(p_1>\tau h_1\) or \(p_1+p_2>\tau h_2\); in either case
all caps vanish.  Thus \(P=H_0\) there, with a uniform
positive lower bound.  We have used the concave limiting
density \(H_0\) throughout; only the model's integrated
areas are approximated by bordered pieces.

Construct the cubic family with the chosen \(d_1,d_2\).
Remove from its base \(C\) the marked points and all
intersections with \(\ell\cup Q\).  The remaining curve
has genus one and admits connected bordered exhaustions
of genus one.  Lemma~\ref{six:toric-model} supplies the
models \eqref{six:surface-form}, with uniform limiting
area \(H_0\), over these exhaustions.
Lemma~\ref{surf:packing} therefore embeds disjoint closed
balls of capacities \(R_i<r_i\), \(1\le i\le k\), compactly
in one such model, with neighborhood embeddings.
The entire chosen model has \(p_1\le h_1<c\), is away
from the double locus, and avoids the inverse images
of \(\ell\cup Q\) under the map to the plane.

Transfer this compact region into a nearby fiber
\(\widehat V\).  Lemma~\ref{six:moment-transfer} keeps it
below \(p_1<c\) and preserves the stated avoidance.
Apply the cut of Lemma~\ref{six:cut}.  The balls now lie
in \(Y\) off infinity and off the two forbidden vertical
submanifolds over \(D\cap E\) and \(D\cap H_\infty\).
The cut form is invariant and is in exactly the
polarization class of Lemma~\ref{six:quadric}.

Use Proposition~\ref{geo:transfer} for the first normal-cone
degeneration.  The allowed relative Moser identification
of Lemma~\ref{geo:relative-moser} preserves infinity and
the two disjoint forbidden vertical submanifolds.
It follows that the compact packing transfers into
\(T\setminus(E\cup H_\infty)\), with a Kähler form in
class \(H-aE\).  Lemma~\ref{geo:shell}, using relative
Moser for the disjoint divisors \(E,H_\infty\), identifies
this complement symplectically with
\[
 \left\{z\in\C^3:
          a<\pi\sum_{j=1}^3|z_j|^2<1\right\}.
\]
Add the central closed ball of capacity \(R_0<a\).
It is disjoint from all these compact shell balls, and
every source lies inside the open target.  Each operation
was defined on a neighborhood of the relevant compact
set; restricting these neighborhoods if necessary gives
the required neighborhood embeddings of the closed balls.
\end{proof}
